\chapter{Lezione 10 - 24/10}


\Esempio{Isomorfismo \(\phi_{\mathcal{B}}\) e Test di Indipendenza}{
    Sia \(V = \mathbb{R}[x]_{\leq 2}\) lo spazio dei polinomi di grado \(\leq 2\).
    Sia \(\mathcal{B} = (1+x, 1-x, x^2)\) una sua base.
    
    L'isomorfismo delle coordinate è \(\phi_{\mathcal{B}} : \mathbb{R}[x]_{\leq 2} \to \mathbb{R}^3\).
    
    Vogliamo trovare la mappatura per un generico polinomio \(p(x) = a_0 + a_1 x + a_2 x^2\).
    Dobbiamo trovare gli unici scalari \(\alpha_1, \alpha_2, \alpha_3\) tali che:
    \[
    a_0 + a_1 x + a_2 x^2 = \alpha_1(1+x) + \alpha_2(1-x) + \alpha_3(x^2)
    \]

    
    Raggruppando per potenze di \(x\):
    \[
    a_0 + a_1 x + a_2 x^2 = (\alpha_1 + \alpha_2) \cdot 1 + (\alpha_1 - \alpha_2) \cdot x + \alpha_3 \cdot x^2
    \]
    
    Uguagliando i coefficienti, otteniamo il sistema:
    \[
    \begin{cases}
    \alpha_1 + \alpha_2 = a_0 \\
    \alpha_1 - \alpha_2 = a_1 \\
    \alpha_3 = a_2
    \end{cases}
    \]
    
    Risolviamo per \(\alpha_1, \alpha_2, \alpha_3\):
    \begin{itemize}
        \item Sommando la prima e la seconda equazione: \(2\alpha_1 = a_0 + a_1 \implies \alpha_1 = \frac{1}{2}a_0 + \frac{1}{2}a_1\)
        \item Sottraendo la seconda dalla prima: \(2\alpha_2 = a_0 - a_1 \implies \alpha_2 = \frac{1}{2}a_0 - \frac{1}{2}a_1\)
        \item Dalla terza equazione: \(\alpha_3 = a_2\)
    \end{itemize}
    
    Quindi, l'isomorfismo \(\phi_{\mathcal{B}}\) mappa il polinomio nelle sue coordinate:
    \[
    \phi_{\mathcal{B}}(a_0 + a_1 x + a_2 x^2) = \left( \frac{1}{2}a_0 + \frac{1}{2}a_1, \, \frac{1}{2}a_0 - \frac{1}{2}a_1, \, a_2 \right)
    \]
    
    ---
    
    \textbf{Applicazione:} Vogliamo testare la lineare indipendenza dell'insieme \(S \subseteq V\):
    \[
    S = \{ p_1, p_2, p_3 \} = \{ x^2+7x-5, \quad 2x-3, \quad 70+5x-3x^2 \}
    \]
    
    
}

\Esempio{}{
    Per l'osservazione "Gli isomorfismi conservano la L.I.", \(S\) è L.I. \(\iff\) l'insieme delle sue immagini \(\phi_{\mathcal{B}}(S)\) è L.I. in \(\mathbb{R}^3\).
    
    Calcoliamo i vettori delle coordinate:
    \begin{itemize}
        \item \( p_1 = -5 + 7x + 1x^2 \implies a_0=-5, a_1=7, a_2=1 \)
        \[ \phi_{\mathcal{B}}(p_1) = \left( \frac{-5+7}{2}, \frac{-5-7}{2}, 1 \right) = (1, -6, 1) \]
        
        \item \( p_2 = -3 + 2x + 0x^2 \implies a_0=-3, a_1=2, a_2=0 \)
        \[ \phi_{\mathcal{B}}(p_2) = \left( \frac{-3+2}{2}, \frac{-3-2}{2}, 0 \right) = \left(-\frac{1}{2}, -\frac{5}{2}, 0\right) \]
        
        \item \( p_3 = 70 + 5x - 3x^2 \implies a_0=70, a_1=5, a_2=-3 \)
        \[ \phi_{\mathcal{B}}(p_3) = \left( \frac{70+5}{2}, \frac{70-5}{2}, -3 \right) = \left(\frac{75}{2}, \frac{65}{2}, -3\right) \]
    \end{itemize}
    
    Per verificare se questi 3 vettori in \(\mathbb{R}^3\) sono L.I., li inseriamo come \textbf{righe} di una matrice \(A\) e ne calcoliamo il rango.
    \[
    A = \begin{pmatrix}
    1 & -6 & 1 \\
    -\frac{1}{2} & -\frac{5}{2} & 0 \\
    \frac{75}{2} & \frac{65}{2} & -3
    \end{pmatrix}
    \]
    L'insieme \(S\) è linearmente indipendente se e solo se \(\dim(L(\phi_{\mathcal{B}}(S))) = 3\), che a sua volta è vero se e solo se \(\operatorname{rango}(A) = 3\).
    
}

\bigskip \bigskip \bigskip

\section{Rango}
\Definizione{Rango di una Matrice}{
    Sia \(A\) una matrice a coefficienti in un campo \(K\).
    
    Il \textbf{rango} di \(A\) (scritto \(\operatorname{rango}(A)\) o \(\operatorname{rk}(A)\)) è un numero intero non negativo che corrisponde a:
    \begin{itemize}
        \item La dimensione dello spazio vettoriale generato dalle \textbf{colonne} della matrice (detto rango per colonne).
        \item La dimensione dello spazio vettoriale generato dalle \textbf{righe} della matrice (detto rango per righe).
    \end{itemize}
    Per il \textbf{Teorema del Rango}, questi due valori (dimensione dello spazio delle righe e dimensione dello spazio delle colonne) sono sempre uguali.
    
}

\bigskip \bigskip \bigskip

\section{Teorema del Rango}
\Teorema{Teorema del Rango}{
    Sia \(A\) una matrice a coefficienti in un campo \(K\).
    Il rango per righe di \(A\) (la dimensione del sottospazio generato dalle righe) è uguale al rango per colonne di \(A\) (la dimensione del sottospazio generato dalle colonne).
    
    Di conseguenza, il rango di una matrice è uguale al rango della sua trasposta:
    \[
    \operatorname{rango}(A) = \operatorname{rango}(A^t)
    \]
}



\BoxBlue{Osservazione}{Spazio delle Righe e Spazio delle Colonne}{
    Prendiamo una matrice \(A \in M_{4,5}(K)\) (una matrice con 4 righe e 5 colonne).
    
    \begin{itemize}
        \item Le sue \textbf{righe} \(a^i\) (che sono 4) sono vettori in \(K^5\).
        Esse generano un sottospazio vettoriale \(W_R = L(a^1, a^2, a^3, a^4) \subseteq K^5\).
        
        \item Le sue \textbf{colonne} \(a_j\) (che sono 5) sono vettori in \(K^4\).
        Esse generano un sottospazio vettoriale \(W_C = L(a_1, a_2, a_3, a_4, a_5) \subseteq K^4\).
    \end{itemize}
    
    Come hai giustamente osservato, \(W_R\) e \(W_C\) sono sottospazi vettoriali \textbf{diversi}, in quanto appartengono a spazi vettoriali "ambiente" diversi (\(K^5\) e \(K^4\)).
    
    Il Teorema del Rango ci dice che, nonostante i due sottospazi siano diversi e vivano in spazi diversi, essi hanno \textbf{la stessa dimensione}:
    \[
    \dim(W_R) = \dim(W_C) = \operatorname{rango}(A)
    \]
}

\Proposizione{Invarianza dello Spazio delle Righe}{
    Sia \(A \in M_{m,n}(K)\) e sia \(B \in M_{m,n}(K)\) una matrice ottenuta da \(A\) mediante un numero finito di operazioni elementari sulle righe.
    
    Siano \(\mathcal{R}(A) = L(a^1, \dots, a^m)\) lo spazio generato dalle righe di \(A\), e \(\mathcal{R}(B) = L(b^1, \dots, b^m)\) lo spazio generato dalle righe di \(B\).
    
    Allora i due spazi delle righe coincidono:
    \[ \mathcal{R}(A) = \mathcal{R}(B) \]
    Di conseguenza, il loro rango (che è la dimensione di questo spazio) è identico:
    \[ \operatorname{rango}(A) = \operatorname{rango}(B) \]
    
    \Dimostrazione{
        L'idea centrale è che ogni \textbf{operazione elementare sulle righe non modifica lo spazio generato dalle righe} di una matrice.\\[6pt]

        \textbf{1.} Mostriamo che \(\mathcal{R}(B) \subseteq \mathcal{R}(A)\):  
        in ciascun tipo di operazione elementare, ogni nuova riga di \(B\) è una combinazione lineare delle righe di \(A\).  
        Poiché lo spazio delle righe è chiuso per combinazioni lineari, segue che \(\mathcal{R}(B) \subseteq \mathcal{R}(A)\).\\[6pt]

        \textbf{2.} Mostriamo che \(\mathcal{R}(A) \subseteq \mathcal{R}(B)\):  
        ogni operazione elementare è \textbf{invertibile} (ha un'operazione inversa dello stesso tipo), quindi possiamo ottenere \(A\) da \(B\).  
        Applicando la stessa logica del punto precedente, si ha che ogni riga di \(A\) è combinazione lineare delle righe di \(B\).\\[6pt]
        
        \textbf{Conclusione:}  
        Dalle due inclusioni \(\mathcal{R}(A) \subseteq \mathcal{R}(B)\) e \(\mathcal{R}(B) \subseteq \mathcal{R}(A)\) segue che \(\mathcal{R}(A) = \mathcal{R}(B)\).  
        Poiché il rango è la dimensione dello spazio delle righe, concludiamo che:
        \[
            \operatorname{rango}(A) = \operatorname{rango}(B).
        \]
    }

}

\section{Rango di matrici a gradini}
\Proposizione{Rango di Matrici a Gradini}{
    Sia \(A \in M_{m,n}(K)\) una matrice ridotta a gradini, allora il rango di A è uguale al numero dei suoi pivot.
    
    \Dimostrazione{
        Sia \(p\) il numero di righe non nulle (e quindi il numero di pivot) della matrice \(A\).
        Le righe non nulle siano \(a^1, \dots, a^p\), mentre le righe \(a^{p+1}, \dots, a^m\) sono tutte uguali al vettore nullo \(\bar{0}\).
        
        Lo spazio delle righe di \(A\) è:
        \[ \mathcal{R}(A) = L(a^1, \dots, a^p, a^{p+1}, \dots, a^m) = L(a^1, \dots, a^p) \]
        (poiché i vettori nulli non contribuiscono allo span).\\[6pt]
        
        \textbf{Tesi:} \(\operatorname{rango}(A) = p\). Dobbiamo quindi mostrare che le righe p non nulle sono \textbf{linearmente indipendenti}.\\[6pt]
        
        Procediamo per induzione su \(p\):
        
        \begin{itemize}
            \item \textbf{Base (p = 0):}
            Se \(A\) è la matrice nulla, non ha righe non nulle, quindi non ha pivot \Rightarrow 
            \(\mathcal{R}(A) = \{\bar{0}\}\) e \(\operatorname{rango}(A) = 0\).
            
            \item \textbf{Passo induttivo (p-1 \(\implies\) p):}
            Supponiamo (Ipotesi Induttiva) che qualsiasi matrice a gradini con \(p-1\) pivot abbia rango \(p-1\).
            
            Sia ora \(A\) una matrice con \(p > 0\) pivot. Visualizziamo la sua struttura:
            \[
            A = \begin{pmatrix}
            0 & \dots & p_1 & \dots & \dots & \dots \\
            0 & \dots & 0 & \dots & p_2 & \dots \\
            \vdots & & \vdots & & \vdots & \\
            0 & \dots & 0 & \dots & 0 & \dots & p_p & \dots \\
            0 & \dots & 0 & \dots & 0 & \dots & 0 & \dots \\
            \vdots & & \vdots & & \vdots & & \vdots &
            \end{pmatrix}
            \begin{matrix}
            \leftarrow a^1 \\
            \leftarrow a^2 \\
            \vdots \\
            \leftarrow a^p \\
            \leftarrow \bar{0} \\
            \vdots
            \end{matrix}
            \]
            
            Sia \(A'\) la matrice ottenuta da \(A\) \textbf{cancellando la prima riga \(a^1\)}.
            La matrice \(A'\) è ancora una matrice a gradini, quindi per ipotesi induttiva le sue righe 
            \(a^2, \dots, a^p\) sono linearmente indipendenti.
            
            Indichiamo con \(j_1\) la colonna del primo pivot \(p_1\) (in \(a^1\)).  
            Per definizione di forma a gradini, in tale colonna tutte le righe sottostanti hanno zero:
            \[
                (a^1)_{j_1} = p_1 \neq 0, 
                \quad (a^k)_{j_1} = 0 \ \text{per } k = 2, \dots, p.
            \]
            Quindi nessuna combinazione lineare di \(a^2, \dots, a^p\) può produrre \(a^1\), 
            poiché nella posizione \(j_1\) tutte avrebbero zero.  
            Segue che \(a^1 \notin L(a^2, \dots, a^p)\).\\[6pt]
            
            \textbf{Conclusione:}
            Per il \textbf{teorema fondamentale sull'indipendenza lineare}, se abbiamo un insieme \textbf{linearmente indipendente} \(S = \{a^2, \dots, a^p\}\)
            e aggiungiamo un vettore \(a^1\) che non appartiene allo span di \(S\), il nuovo insieme \(\{a^1\} \cup S = \{a^1, a^2, \dots, a^p\}\) \textbf{è ancora linearmente indipendente}.\\[4pt]
            
            Dato che lo spazio delle righe \(\mathcal{R}(A)\) è generato da queste \(p\) righe linearmente indipendenti, la sua dimensione è \(p\).
            \[
            \operatorname{rango}(A) = \dim(\mathcal{R}(A)) = p
            \]
        \end{itemize}
    }
}


\BoxBlue{Osservazione}{(Definizione Equivalente di Rango)}{
    Sia \(A \in M_{m,n}(K)\).  
    Il \textbf{rango} di \(A\), è la dimensione dello spazio vettoriale generato dalle sue righe.  
    Per il \textit{Teorema del Rango}, questa dimensione coincide con quella dello spazio generato dalle colonne\\[4pt]
    
    Equivalentemente, il rango può essere visto come:
    \begin{itemize}
        \item il massimo numero di \textbf{righe linearmente indipendenti};
        \item oppure, in modo del tutto analogo, il massimo numero di \textbf{colonne linearmente indipendenti}.
    \end{itemize}
}

\section{Teorema di Rouché-Capelli}

\Teorema{di Rouché-Capelli}{
    Sia \(\Sigma\) un sistema lineare \(Ax = b\), con \(A \in M_{m,n}(K)\) (m equazioni in n incognite), e sia 
    \(C = (A \mid b)\) la matrice completa associata al sistema.
    
    Allora, \(\Sigma\) è \textbf{compatibile} se e solo se:
    \[
        \operatorname{rango}(A) = \operatorname{rango}(C)
    \]

    \Dimostrazione{
        Applichiamo l'algoritmo di Gauss a \(C\) per ottenere la forma ridotta a gradini:
        \[
            C' = (A' \mid b').
        \]
        
        Ricordiamo due fatti fondamentali:
        \begin{enumerate}
            \item Le operazioni elementari di riga non alterano il rango:
            \[
                \operatorname{rango}(A) = \operatorname{rango}(A'), \quad
                \operatorname{rango}(C) = \operatorname{rango}(C').
            \]
            \item Il sistema originale \(Ax=b\) ha le stesse soluzioni del sistema ridotto \(A'x = b'\).\\[4pt]
        \end{enumerate}
        
        \textbf{Analisi delle righe di tipo impossibile:}  
        Il sistema è \textbf{incompatibile} se e solo se in \(C'\) esiste una riga del tipo:
        \[
            0x_1 + 0x_2 + \dots + 0x_n = \beta_i, \quad \beta_i \neq 0
        \]
        cioè una riga in cui la parte di \(A'\) è nulla ma il termine noto non lo è.  
        Tale riga possiede un pivot nell'ultima colonna (colonna dei termini noti).\\[4pt]
        
        \textbf{Relazione tra rango e compatibilità:}
        
        \begin{itemize}
            \item \textbf{Caso compatibile:}  
            Nessuna riga di \(C'\) è del tipo \((0, \dots, 0 \mid \beta_i \neq 0)\).  
            \[
                C' = 
                \begin{pmatrix}
                1 & 2 & 0 & | & 5 \\
                0 & 0 & 1 & | & 3 \\
                0 & 0 & 0 & | & 0
                \end{pmatrix}
                \quad \Rightarrow \quad
                \operatorname{rango}(A') = \operatorname{rango}(C') = 2
            \]
            Tutti i pivot si trovano nella parte di \(A'\), quindi il sistema è compatibile.
            
            \item \textbf{Caso incompatibile:}  
            Esiste almeno una riga del tipo \((0, \dots, 0 \mid \beta_i \neq 0)\), che crea un'equazione impossibile.  
            \[
                C' = 
                \begin{pmatrix}
                1 & 2 & 0 & | & 5 \\
                0 & 0 & 1 & | & 3 \\
                0 & 0 & 0 & | & 4
                \end{pmatrix}
                \quad \Rightarrow \quad
                \operatorname{rango}(C') = 3 > \operatorname{rango}(A') = 2
            \]
            Il pivot nell'ultima colonna indica che il sistema non ha soluzioni.
        \end{itemize}
        
        Abbiamo così dimostrato entrambe le implicazioni, completando la dimostrazione.
        \[
            \Sigma \text{ compatibile } \iff \operatorname{rango}(A) = \operatorname{rango}(C)
        \]
    }
}



\Teorema{Rappresentazione Cartesiana di un Sottospazio}{
    Sia \(W \subseteq K^n\) un sottospazio vettoriale con \(\dim(W) = h\).
    
    Allora esiste un sistema lineare omogeneo \(\Sigma_0: Ax = \bar{0}\) composto da \(n - h\) equazioni in \(n\) incognite, con \(\operatorname{rango}(A) = n-h\), tale che l'insieme delle sue soluzioni \(S_0\) coincida con \(W\).
    \[ S_0 = W \]

    \Dimostrazione{
    Sia \(\mathcal{B}_W = \{w_1, \dots, w_h\}\) una base di \(W\).
    
    Scriviamo le componenti dei vettori della base (rispetto alla base canonica di \(K^n\)) come colonne:
    \[
    w_1 = \begin{pmatrix} a_{11} \\ a_{21} \\ \vdots \\ a_{n1} \end{pmatrix}, \quad
    w_2 = \begin{pmatrix} a_{12} \\ a_{22} \\ \vdots \\ a_{n2} \end{pmatrix}, \quad \dots, \quad
    w_h = \begin{pmatrix} a_{1h} \\ a_{2h} \\ \vdots \\ a_{nh} \end{pmatrix}
    \]
    Sia \(A = (w_1 | \dots | w_h)\) la matrice \(n \times h\) che ha questi vettori come colonne. Poiché sono una base, \(\operatorname{rango}(A) = h\).
    
    Un generico vettore \(x = (x_1, \dots, x_n)^t\) appartiene a \(W\) se e solo se è combinazione lineare di \(w_1, \dots, w_h\).
    
    Come hai scritto tu, costruiamo la matrice \(C\) (orlata) accostando \(x\) alle colonne di \(A\):
    \[
    C = (A | x) = \begin{pmatrix}
    a_{11} & \dots & a_{1h} & x_1 \\
    a_{21} & \dots & a_{2h} & x_2 \\
    \vdots & \ddots & \vdots & \vdots \\
    a_{n1} & \dots & a_{nh} & x_n
    \end{pmatrix}
    \]
    Per il Teorema di Rouché-Capelli (visto come test di appartenenza), \(x \in W\) se e solo se:
    \[
    \operatorname{rango}(C) = \operatorname{rango}(A)
    \]
    Poiché \(\operatorname{rango}(A) = h\), la condizione è:
    \[
    x \in W \iff \operatorname{rango}(C) = h
    \]
}


}

